\paragraph{Problem F.2.4 (Open in this incomplete return).}
The simultaneous integral conjugacy and affine-twist
classification is not complete.
The existence alternative in F.4 is proved by Theorem F.4.1.
We have reduced the cusp
to a single matrix with its exact centralizer, bounded every
possible finite linear order, run the complete 38-type and
19-type characteristic-polynomial filters, and excluded all
self-dual candidates without a finite search.  The smallest
remaining $k=2$ pattern, with one global covector, no global vector
and local fixed ranks $(2,2,3)$, is realized integrally by
Theorem F.2.5, including its primitive cusp and free affine actions.
Its integral total torsion is settled in F.2.6, and
Theorems F.3.5--F.3.7 construct its global holomorphic periods,
finite affine quotients and compact cusp filling.

The attempted finite-order enumeration does not bound the
relative positions of the finite matrices under the cusp
centralizer.  The local filling counterexample in F.2.3
prevents using those matrices alone as an integral homology
algorithm.  The parity proof fails for a non-self-dual module
exactly because $b$ can be even while $a=1$; (F.4) records all
such first-cohomology possibilities.  A polarization on a
three-torus period family would make its representation
self-dual and is therefore ruled out, but the question permits
unpolarized families.  The proposed algebraic abelian-surface
source is nevertheless excluded by Theorem F.3.4 through its
filtration and semisimple middle quotient.
For the explicit candidate, the order-five filling has
$H_1=\mathbb Z^2$, while the order-four filling has
$H_1=\mathbb Z^2\oplus\mathbb Z/2$.
Its low-degree total torsion issue was genuinely two-primary.
The finite covering images in degrees two and three have index
one; the integral parabolic groups in coefficient degrees two
and four vanish, and that in degree three is free of rank three.
These assertions are checked by suspension periods, middle
pairings and local restriction Smith forms in the same code.
The nonzero torsion transgressions in F.2.6 settle that issue.
Its exact smooth type is determined in Theorem F.4.1 by the
characteristic-vector orbit, Wall's normal invariant and the
published inertia-group calculation.  The remaining work in F.2
is an integral classification of all relative positions and twists,
with the Kuhn filling fixed.

